\documentclass[11pt,a4paper]{article}
% Shared preamble for the three Phase 2 deliverables
\usepackage[utf8]{inputenc}
\usepackage[T1]{fontenc}
\usepackage{lmodern}
\usepackage[margin=2.3cm]{geometry}
\usepackage{graphicx}
\usepackage{xcolor}
\usepackage{booktabs}
\usepackage{enumitem}
\usepackage{titlesec}
\usepackage{fancyhdr}
\usepackage{tcolorbox}
\usepackage{float}
\usepackage{caption}
\usepackage{parskip}
\usepackage{listings}
\usepackage{array}
\usepackage{hyperref}

\definecolor{primary}{RGB}{44,62,80}
\definecolor{accent}{RGB}{52,152,219}
\definecolor{insightbg}{RGB}{235,245,255}
\definecolor{bodytext}{RGB}{50,50,50}
\definecolor{kw}{RGB}{0,92,170}
\definecolor{cm}{RGB}{110,110,110}
\definecolor{st}{RGB}{163,21,21}
\definecolor{codebg}{RGB}{248,249,251}

\hypersetup{colorlinks=true, linkcolor=accent, urlcolor=accent,
  pdfauthor={Saanvi Grover, Aditya Sharma}}

\titleformat{\section}{\Large\bfseries\color{primary}}{}{0em}{}[\color{accent}\titlerule]
\titleformat{\subsection}{\large\bfseries\color{primary}}{}{0em}{}
\titleformat{\subsubsection}{\normalsize\bfseries\color{primary}}{}{0em}{}

\pagestyle{fancy}
\fancyhf{}
\fancyfoot[C]{\thepage}
\renewcommand{\headrulewidth}{0pt}

\lstdefinelanguage{SQLplus}[]{SQL}{
  morekeywords={WITH, OVER, PARTITION, RANK, NTILE, PERCENT_RANK, COALESCE, ROUND,
                CROSS, RETURNING, WINDOW, ROWS, PRECEDING, FOLLOWING, CURRENT, UNBOUNDED},
  sensitive=false
}
\lstset{
  language=SQLplus,
  basicstyle=\ttfamily\footnotesize\color{bodytext},
  keywordstyle=\color{kw}\bfseries,
  commentstyle=\color{cm}\itshape,
  stringstyle=\color{st},
  backgroundcolor=\color{codebg},
  frame=single, rulecolor=\color{accent}, framerule=0.4pt,
  numbers=left, numberstyle=\tiny\color{cm}, numbersep=8pt,
  breaklines=true, showstringspaces=false, tabsize=4,
  xleftmargin=1.6em, aboveskip=6pt, belowskip=6pt, columns=fullflexible, keepspaces=true
}

\newtcolorbox{insightbox}[1][KEY INSIGHT]{
  colback=insightbg, colframe=accent, boxrule=0.5pt, arc=2pt,
  left=8pt, right=8pt, top=6pt, bottom=6pt,
  fonttitle=\bfseries\color{accent}\small, title=#1
}

\newcommand{\entropytitle}[2]{%
\begin{titlepage}
\centering
\vspace*{3.5cm}
{\Huge\bfseries\color{primary} #1\par}
\vspace{0.6cm}
{\Large\color{gray} #2\par}
\vspace{1.2cm}
{\color{accent}\rule{8cm}{1pt}\par}
\vspace{1.2cm}
{\large\color{gray} Round 1, Phase 2: Analytical Core\par}
\vspace{1.5cm}
{\large
\begin{tabular}{rl}
\textbf{Competition:} & Data Vortex $|$ AARUUSH'26 \\
\textbf{Theme:} & Rebuilding the Social Engine \\
\textbf{Team:} & Entropy \\
\textbf{Members:} & Saanvi Grover \& Aditya Sharma \\[0.4cm]
\textbf{Questions:} & E3 $\cdot$ M4 $\cdot$ H4 \\
\textbf{Database:} & SQLite (\texttt{Entropy\_social\_engine.db}) \\
\end{tabular}
\par}
\vfill
\end{titlepage}
}

\hypersetup{pdftitle={Phase 2 Logic Explanation -- Team Entropy}}

\begin{document}
\color{bodytext}
\entropytitle{Logic Explanation}{Approach, design decisions and validation for E3, M4 and H4}

\section{Why These Three Questions}

We chose one question per level that (a) demands a genuinely different SQL technique, (b) forces an explicit decision about the NULLs left by the corruption, and (c) can be checked against a finding from our Phase 1 EDA rather than only returning rows.

\begin{table}[H]
\centering\small
\begin{tabular}{@{}llp{7.2cm}@{}}
\toprule
\textbf{Question} & \textbf{Core technique} & \textbf{Phase 1 finding it tests} \\
\midrule
E3 Average engagement by platform & \texttt{GROUP BY} + \texttt{RANK()} + cross-joined benchmark & Platform choice does not affect engagement (Kruskal--Wallis $p=0.37$) \\
M4 Platform behaviour, $\geq$30k followers & Join + filter + two aggregated CTEs joined on platform & Follower count has no predictive power (Spearman $\rho = 0.025$) \\
H4 Follower-to-engagement anomaly & Three-level CTE chain + \texttt{NTILE(10)} + \texttt{PERCENT\_RANK()} & Engagement is independent of user attributes; ``anomalies'' should occur at chance rate \\
\bottomrule
\end{tabular}
\end{table}

\section{Shared Design Decisions}

\subsection{Handling the corrupted NULLs}

The cleaned dataset intentionally keeps NULLs in \texttt{platform} (1,784 posts), \texttt{likes} (1,814) and \texttt{text\_content}. Two different rules are used, and the choice is deliberate:

\begin{itemize}[leftmargin=*]
  \item \textbf{Averages (E3, M4): exclude the row.} \texttt{AVG()} skips NULLs on its own, but \texttt{AVG(likes + shares + comments)} would silently drop a row from the total while still counting its shares in \texttt{AVG(shares)}. Filtering \texttt{likes IS NOT NULL} once in the \texttt{WHERE} clause keeps all four averages on the \emph{same} denominator. Substituting 0 would bias every average downward by about 15\%.
  \item \textbf{Sums per user (H4): \texttt{COALESCE(likes, 0)}.} A user's total is a sum across many posts. Dropping a post because its like count is corrupted would also discard its (valid) shares and comments, penalising users at random. Treating the missing likes as 0 keeps every post and the bias is uniform across users because the corruption is MCAR.
\end{itemize}

\subsection{Defining ``engagement''}

Everywhere, engagement $=$ likes $+$ shares $+$ comments, as defined by question E2 in the brief. E3 also reports the three components separately because the question asks for them.

\subsection{Ranking inside SQL}

Rankings are produced with \texttt{RANK() OVER (ORDER BY ...)} rather than relying on \texttt{ORDER BY} alone, so the rank survives if the consumer re-sorts the result and ties are made visible.

\section{E3 -- Average Engagement by Platform}

\subsection{Approach}
\begin{enumerate}[leftmargin=*]
  \item \textbf{\texttt{platform\_stats} CTE}: filter to rows with a known platform and a non-NULL like count, then \texttt{GROUP BY platform} to compute post count and the four averages.
  \item \textbf{\texttt{overall} CTE}: the grand mean total engagement over the same filtered rows. It is a single scalar, so a \texttt{CROSS JOIN} attaches it to every platform row.
  \item \textbf{Final \texttt{SELECT}}: \texttt{RANK()} by average total engagement, plus \texttt{pct\_vs\_overall} = percentage deviation from the grand mean.
\end{enumerate}

\subsection{Why add the benchmark column}
Returning only the ranking answers ``which platform is highest'' but hides \emph{by how much}. The extra column shows Instagram is first by just $+0.78\%$ and Twitter last at $-1.42\%$, which is the real answer to the question.

\subsection{Validation}
\begin{itemize}[leftmargin=*]
  \item Sum of \texttt{post\_count} = 8,662 = 12,000 $-$ 1,784 (NULL platform) $-$ 1,814 (NULL likes) $+$ 260 (both NULL), matching pandas.
  \item \texttt{avg\_total\_engagement} equals \texttt{avg\_likes + avg\_shares + avg\_comments} for every row (linearity of the mean over the same denominator).
  \item Kruskal--Wallis across the five platforms on these rows: $H = 2.97$, $p = 0.56$.
\end{itemize}

\section{M4 -- Platform Behaviour by High-Follower Users}

\subsection{Approach}
\begin{enumerate}[leftmargin=*]
  \item \textbf{\texttt{high\_follower\_posts} CTE}: join \texttt{posts} to \texttt{users} on \texttt{user\_id} (indexed on both sides), keep users with \texttt{follower\_count >= 30000}, and compute engagement per post. The same NULL rule as E3 applies.
  \item \textbf{\texttt{cohort\_by\_platform} CTE}: aggregate per platform: distinct users, posts and average engagement.
  \item \textbf{\texttt{baseline\_by\_platform} CTE}: the same per-platform average across \emph{all} users.
  \item \textbf{Final \texttt{SELECT}}: join cohort to baseline on \texttt{platform}, rank by the cohort average and report \texttt{cohort\_minus\_baseline}.
\end{enumerate}

\subsection{Why compare against a baseline}
``Which platform gives high-follower users the most engagement'' only matters if high-follower users behave differently from everyone else. Without the baseline, Instagram at 4,145 looks like a finding; with it, the gap to the all-user Instagram average is $+105$ ($\approx 2.6\%$), and three of the five platforms are \emph{below} their baseline.

\subsection{Validation}
\begin{itemize}[leftmargin=*]
  \item 594 users have $\geq$ 30,000 followers; the five \texttt{high\_follower\_users} counts (394--419) overlap because users post on several platforms.
  \item Kruskal--Wallis across platforms within the cohort: $H = 6.22$, $p = 0.18$.
  \item Mann--Whitney U, cohort vs.\ all other users' posts: $p = 0.98$; the cohort is indistinguishable from the rest.
\end{itemize}

\section{H4 -- Follower-to-Engagement Anomaly}

\subsection{Approach: three levels of analysis}
\begin{enumerate}[leftmargin=*]
  \item \textbf{Post $\rightarrow$ user (\texttt{user\_engagement})}: join and \texttt{GROUP BY} user to obtain \texttt{post\_count} and \texttt{total\_engagement}. Every one of the 1,500 users has at least one post, so an inner join loses nobody.
  \item \textbf{User $\rightarrow$ population (\texttt{ranked\_users})}: \texttt{NTILE(10) OVER (ORDER BY total\_engagement DESC)} assigns each user a decile; decile 1 is exactly the top 10\% (150 users). \texttt{PERCENT\_RANK()} gives a finer percentile for reporting, and \texttt{COUNT(*) OVER ()} carries the population size without another subquery.
  \item \textbf{Population $\rightarrow$ anomalies}: filter \texttt{engagement\_decile = 1 AND follower\_count < 5000}, then \texttt{RANK()} the survivors and derive \texttt{avg\_engagement\_per\_post}.
\end{enumerate}

\subsection{Why \texttt{NTILE} rather than a hard-coded cut-off}
A threshold such as \texttt{total\_engagement > 43948} would be correct today and wrong the moment a row changes. \texttt{NTILE(10)} defines ``top 10\%'' relative to the data, so the query remains valid on any refresh of the table. \texttt{PERCENT\_RANK} is included alongside it because \texttt{NTILE} alone cannot show \emph{how far} inside the top decile each user sits.

\subsection{Why report \texttt{avg\_engagement\_per\_post}}
The question is phrased around \emph{total} engagement, but total $=$ posts $\times$ average. Reporting both lets the reader see immediately whether a user is in the top decile because their posts perform unusually well or simply because they post more. All 16 anomalies average 2,939--4,849 per post, i.e.\ around the platform-wide mean of $\sim$4,000, while posting 10--18 times against a population mean of 8.0.

\subsection{Validation}
\begin{itemize}[leftmargin=*]
  \item Decile 1 contains exactly 150 users (1,500 / 10).
  \item 131 users (8.73\%) have $<$ 5,000 followers. Under independence the expected number in the top decile is $150 \times 0.0873 = 13.1$; 16 are observed (binomial $p = 0.39$).
  \item Spearman correlation, post count vs.\ total engagement: $\rho = 0.89$; follower count vs.\ total engagement: $\rho = -0.03$ ($p = 0.30$).
\end{itemize}

\section{Tooling}

SQLite 3 via Python \texttt{sqlite3} and pandas inside \texttt{Entropy\_01\_SQL\_Analysis.ipynb}; statistical checks with \texttt{scipy.stats}; screenshots and charts rendered with matplotlib by \texttt{Entropy\_run\_queries.py} from the exact query result sets.

\end{document}
